\chapter{Propriétés des facteurs d'homothétie}\label{annex:homo}
\minitoc

\section{Loi de puissance des facteurs homothétiques}

\begin{Theo}[Loi de puissance]
Si le facteur
\(
f(\lambda)
\)
est bien défini indépendamment de \(n>0\) (ce qui est le cas pour les solutions homothétiques),
alors \(f\) est une loi de puissance.
\end{Theo}

\begin{proof}
Posons \(g(n) = \mu'(n)>0\) ou \(<0\) (\ie $\mu$ strictement monotone). La définition de \(f\) équivaut à l’existence d’une fonction
\(\chi(\lambda)=1/f(\lambda)\) telle que
\[
g(\lambda n)=\chi(\lambda)\,g(n)\qquad(\forall\,\lambda,n>0).
\]

En prenant \(n=1$, on a \(\chi(\lambda)=g(\lambda)/g(1)\). Donc, pour tous \(a,b>0\),
\[
\chi(ab)=\frac{g(ab)}{g(1)}=\frac{\chi(a)\,g(b)}{g(1)}=g(a)\,g(b),
\]
c’est-à-dire que \(\chi\) est \emph{multiplicative}. Sous une hypothèse physique très faible
(continuité, mesurabilité ou simple localement bornée), toute fonction multiplicative sur
\(\mathbb{R}_+^\ast\) est de la forme
\[
\chi(\lambda)=\lambda^{\alpha-1}
\quad\Rightarrow\quad
f(\lambda)=\lambda^{1-\alpha}.
\]
\end{proof}

\section{Équivalence entre $f(\lambda)$ et $\mu(n)$}

\begin{Theo}
Les expressions
\[
f(\lambda) = \lambda^{1-\alpha} \quad \text{et} \quad \mu (n) \propto n^\alpha
\]
sont équivalentes.
\end{Theo}

\begin{proof}

\begin{enumerate}
	\item 1. Si \(\mu(n)=C\,n^\alpha\) avec \(C\neq0\) :

Alors \(\mu'(n)=C\alpha\,n^{\alpha-1}\). Par conséquent,
\[
f(\lambda)=\frac{C\alpha\,n^{\alpha-1}}{C\alpha\,(\lambda n)^{\alpha-1}} = \lambda^{1-\alpha}.
\]

	\item Réciproque : si \(f(\lambda)=\lambda^{1-\alpha}\) pour tout \(\lambda>0\) et \(n>0\) :
Posons \(g(n)=\mu'(n)\). L'hypothèse s'écrit
\[
\frac{g(n)}{g(\lambda n)}=\lambda^{1-\alpha} \quad \Longleftrightarrow \quad g(\lambda n)=\lambda^{\alpha-1} g(n),
\]
pour tout \(n>0\) et tout \(\lambda>0$.

Fixons \(n_0>0\) et définissons \(\varphi(\lambda)\equiv g(\lambda n_0)\). La relation ci-dessus donne
\[
\varphi(\lambda)=\lambda^{\alpha-1}\,\varphi(1),
\]
donc \(\varphi(\lambda)=C_1 \lambda^{\alpha-1}\) avec \(C_1 = g(n_0)\). En remplaçant \(\lambda=x/n_0\), on obtient pour tout \(x>0\)
\[
g(x)=C_1\,x^{\alpha-1}.
\]

Ainsi, \(g(n)=\mu'(n)=C\,n^{\alpha-1}\) avec \(C\) constant. En intégrant (si \(\alpha\neq 0\)), on a
\[
\mu(n) \propto n^\alpha.
\]
Pour \(\alpha=0\), \(\mu'(n)=C\,n^{-1}\) et \(\mu(n)=C\ln n+\text{const}\).
\end{enumerate}
\end{proof}

\begin{Rema}
La démonstration utilise la propriété fonctionnelle multiplicative
\(g(\lambda n)=\lambda^{\alpha-1} g(n)\). Sous une hypothèse faible de continuité (ou dérivabilité) en \(n$, cette équation force la forme de puissance \(g(n)\propto n^{\alpha-1}\). Sans régularité, des solutions pathologiques peuvent exister mais ne sont pas physiquement pertinentes dans le contexte thermodynamique.
\end{Rema}


%\paragraph{Proposition.}
%Si le facteur
%\(
%f(\lambda)
%\)
%est bien défini indépendamment de \(n>0\) (ce qui est le cas pour les solutions homothétiques),
%alors \(f\) est une loi de puissance.
%
%\paragraph{Preuve.}
%Posons \(g(n) = \mu'(n)>0\) ou \(<0\) (\ie $\mu$ strictement monotone). La définition de \(f\) équivaut à l’existence d’une fonction
%\(\chi(\lambda)=1/f(\lambda)\) telle que
%\[
%g(\lambda n)=\chi(\lambda)\,g(n)\qquad(\forall\,\lambda,n>0).
%\]
%En prenant \(n=1\), on a \(\chi(\lambda)=g(\lambda)/g(1)\).
%Donc, pour tous \(a,b>0\),
%\[
%\chi(ab)=\frac{g(ab)}{g(1)}=\frac{\chi(a)\,g(b)}{g(1)}=g(a)\,g(b),
%\]
%c’est-à-dire que \(\chi\) est \emph{multiplicative}. Sous une hypothèse physique très faible
%(continuité/mesurabilité ou simple localement bornée), toute fonction multiplicative sur
%\(\mathbb{R}_+^\ast\) est de la forme
%\[
%\chi(\lambda)=\lambda^{\alpha-1}
%\quad\Rightarrow\quad
%f(\lambda)=\lambda^{1-\alpha}.
%\]
%
%\qed
%
%
%% --- Démonstration que f(λ)=λ^{1-\alpha} et réciproque ---
%\paragraph{Proposition.}
%$f(\lambda) = \lambda^{1-\alpha}$ et $\mu (n) \propto n^\alpha $ sont equivalents.
%
%
%\paragraph{1. Si \(\mu(n)=C\,n^\alpha\) (avec \(C\neq0\)) :}
%Alors \(\mu'(n)=C\alpha\,n^{\alpha-1}\). Par conséquent
%\[
%f(\lambda)=\frac{C\alpha\,n^{\alpha-1}}{C\alpha\,(\lambda n)^{\alpha-1}}
%=\lambda^{1-\alpha}.
%\]
%
%\paragraph{2. Réciproque : si \(f(\lambda)=\lambda^{1-\alpha}\) pour tout \(\lambda>0\) (et tout \(n>0\)) :}
%Posons \(g(n)=\mu'(n)\). L'hypothèse s'écrit
%\[
%\frac{g(n)}{g(\lambda n)}=\lambda^{1-\alpha}
%\quad\Longleftrightarrow\quad
%g(\lambda n)=\lambda^{\alpha-1}\,g(n),
%\]
%pour tout \(n>0\) et tout \(\lambda>0\).
%
%Fixons \(n_0>0\) et définissons \(\varphi(\lambda)\equiv g(\lambda n_0)\). La relation ci-dessus donne
%\[
%\varphi(\lambda)=\lambda^{\alpha-1}\,\varphi(1).
%\]
%Autrement dit \(\varphi(\lambda)=C_1\,\lambda^{\alpha-1}\) pour une constante \(C_1=\varphi(1)=g(n_0)\). En remplaçant \(\lambda=x/n_0\) on obtient pour tout \(x>0\)
%\[
%g(x)=C_1\,x^{\alpha-1}.
%\]
%Ainsi \(g(n)=\mu'(n)=C\,n^{\alpha-1}\) avec \(C\) constant.
%
%En intégrant (en supposant \(\alpha\neq 0\)), 
%\(\mu(n)\propto n^\alpha\). (Pour \(\alpha=0\) on obtient \(\mu'(n)=C\,n^{-1}\) et \(\mu(n)=C\ln n+\text{const}\).)
%
%\paragraph{Remarque sur les hypothèses.}
%La démonstration utilise la propriété fonctionnelle multiplicative
%\(g(\lambda n)=\lambda^{\alpha-1}g(n)\). Sous une hypothèse faible de continuité (ou dérivabilité) en \(n\) cette équation force la forme de puissance \(g(n)\propto n^{\alpha-1}\). Sans régularité, des solutions pathologiques peuvent exister mais ne sont pas physiquement pertinentes dans le contexte thermodynamique.
%
%\qed